The status of the spiral nebulae was first contested in open argument and only afterwards resolved by measurement. The Great Debate was held on 26 April 1920 at the Smithsonian Museum of Natural History in Washington, D.C. Harlow Shapley argued that the spiral nebulae were gaseous clouds lying inside a single large Milky Way system, whereas Heber Curtis argued that the spiral nebulae were distant stellar systems comparable in size to the Milky Way. The issue therefore turned on distance, and distance could be extracted only from objects of known intrinsic brightness.

\slotmarker{S1} By contrast, S Andromedae, which appeared in 1885, was a supernova far more luminous than an ordinary nova. The two classes of outburst were not yet separated in the photometry then available, so the eruptions recorded in the spirals could not yet be made to yield a secure scale.

Shapley's case for a single, all-encompassing system rested on evidence of a quite different kind, namely apparent motion within the spirals themselves, since a nebula visibly turning on a human timescale would have to be small and close at hand. \slotmarker{S2} A later reexamination of van Maanen's plates concluded that the reported rotations of spiral nebulae were measurement artifacts, and that verdict cannot be squared with any claim that subsequent comparisons of photographic plates had confirmed those rotations.

Resolution came from a calibrated distance indicator rather than from debate. Earlier, Henrietta Leavitt established the period-luminosity relation from Cepheid variables in the Magellanic Clouds; the search for such stars in the spirals came only afterwards. \slotmarker{S3} A figure of that magnitude demonstrates that M31 lies outside the boundaries of the Milky Way system. The two placements exclude one another: either M31 lies beyond the Galaxy's confines, or M31 lies within the boundaries of the Milky Way system, and exactly one of these must hold. On the Cepheid evidence, the spirals are stellar systems in their own right.
